
\documentclass{article}
\usepackage{colortbl}
\usepackage{makecell}
\usepackage{multirow}
\usepackage{supertabular}

\begin{document}

\newcounter{utterance}

\centering \large Interaction Transcript for game `taboo', experiment `high\_en', episode 16 with reka{-}flash{-}3.1{-}t0.0.
\vspace{24pt}

{ \footnotesize  \setcounter{utterance}{1}
\setlength{\tabcolsep}{0pt}
\begin{supertabular}{c@{$\;$}|p{.15\linewidth}@{}p{.15\linewidth}p{.15\linewidth}p{.15\linewidth}p{.15\linewidth}p{.15\linewidth}}
    \# & $\;$A & \multicolumn{4}{c}{Game Master} & $\;\:$B\\
    \hline

    \theutterance \stepcounter{utterance}  
    & & \multicolumn{4}{p{0.6\linewidth}}{
        \cellcolor[rgb]{0.9,0.9,0.9}{
            \makecell[{{p{\linewidth}}}]{
                \texttt{\tiny{[P1$\langle$GM]}}
                \texttt{You are playing a collaborative word guessing game in which you have to describe a target word for another player to guess.} \\
\\ 
\texttt{Rules:} \\
\texttt{(a) You have to reply in the form: CLUE: <some text>. Guesses from the other player will start with GUESS.} \\
\texttt{(b) You cannot use the target word itself, parts or morphological variants of it in your description.} \\
\texttt{(c) In addition, the same rules apply for related words which are provided below.} \\
\\ 
\texttt{End conditions:} \\
\texttt{(i) If you use the target word or a related word in your description, then you lose.} \\
\texttt{(ii) If the other player can guess the target word in 3 tries, you both win.} \\
\\ 
\texttt{Let us start.} \\
\\ 
\texttt{This is the target word that you need to describe and that the other player needs to guess:} \\
\\ 
\texttt{field} \\
\\ 
\texttt{Related words are:} \\
\\ 
\texttt{{-} farm} \\
\texttt{{-} grass} \\
\texttt{{-} area} \\
\\ 
\texttt{Important: You are under time pressure, give short descriptions that are to the point!} \\
            }
        }
    }
    & \\ \\

    \theutterance \stepcounter{utterance}  
    & \multicolumn{4}{p{0.6\linewidth}}{
        \cellcolor[rgb]{0.9,0.9,0.9}{
            \makecell[{{p{\linewidth}}}]{
                \texttt{\tiny{[P1$\rangle$GM]}}
                \texttt{<reasoning>} \\
\texttt{I need to come up with a clue for the word "field" without using the word itself, "farm," "grass," or "area." Let me think about what a field is. A field is an open piece of land, usually used for agriculture. But I can't mention farm. Maybe focus on the activities associated with fields. Like planting crops, maybe? But "crops" might be related to farming. Alternatively, think about the environment. Fields are often green, but grass is a related word, so I can't use that either. Maybe something about the sky? Fields are under the sky, but that's too vague. What about the shape? Fields are usually rectangular or large tracts. Or maybe the animals that graze there, like cows or sheep. But "livestock" might be okay. Or maybe the sound of a field, like rustling leaves or wind. "Whispering winds" could work. Or think about sports. A field is where sports are played, like soccer or baseball. But "sports" might be too broad. Alternatively, "where crops grow" but "crops" might be okay. Let me check the related words again: farm, grass, area. So "crops" might be allowed. But maybe "crops" is part of a farm. Hmm. Alternatively, "open space for activities." But "activities" is too vague. Maybe "where tractors drive." But "tractors" relate to farming. Maybe "a place with tall plants swaying." But "tall plants" could be grass, which is a related word. So that's not allowed. Maybe "a large, flat expanse." But "expanse" might be too similar to "area," which is a related word. Wait, "area" is a related word, so I can't use that. So "large, flat space" might be okay. Or "where harvest happens." Harvest is part of farming, so maybe not. Alternatively, "used for growing things." But "growing things" is vague. Maybe "a place where you see crops." But again, "crops" might be related. Alternatively, "a patch of land with scattered plants." Not sure. Maybe think of the four seasons. Fields change with seasons. But how to phrase that without using forbidden words. "Where seasons change the landscape." Maybe. Or "a place where} \\
            }
        }
    }
    & & \\ \\

    \theutterance \stepcounter{utterance}  
    & & & \multicolumn{2}{p{0.3\linewidth}}{
        \cellcolor[rgb]{0.9,0.9,0.9}{
            \makecell[{{p{\linewidth}}}]{
                \texttt{\tiny{[GM$|$GM]}}
                \texttt{abort game} \\
            }
        }
    }
    & & \\ \\

\end{supertabular}
}

\end{document}
